\subsection{Properties of the vectorial integral}

PROPOSITION 1. --- The mapping

\[
(\mathbf {f}, \mu) \mapsto \int \mathbf {f} d \mu
\]

of \(\widetilde{\mathcal{K}} (\mathrm{X};\mathrm{E})\times \mathcal{M}(\mathrm{X};\mathbf{C})\) into \(\mathrm{E}^{\prime *}\) is bilinear.

The proposition follows immediately from Def. 1 of No.~1.

Let \(\mathbf{u}\) be a continuous linear mapping of \(\mathrm{E}\) into a locally convex space \(\mathrm{F}\) ; we know that the transpose \(^t\mathbf{u}\) is a linear mapping of the dual \(\mathrm{F}'\) of \(\mathrm{F}\) into the dual \(\mathrm{E}'\) of \(\mathrm{E}\) ; we shall denote by \(^{tt}\mathbf{u}\) the mapping \(\mathrm{E}^{\prime *} \to \mathrm{F}^{\prime *}\) , the transpose of \(^t\mathbf{u}\) (in the algebraic sense); when \(\mathrm{E}\) and \(\mathrm{F}\) are Hausdorff and are canonically identified with subspaces of \(\mathrm{E}^{\prime *}\) and \(\mathrm{F}^{\prime *}\) , respectively, \(^{tt}\mathbf{u}\) extends the mapping \(\mathbf{u}\) . With these notations:

PROPOSITION 2. --- Let u be a continuous linear mapping of E into a locally convex space F; for every function \(\mathbf{f} \in \widetilde{\mathcal{K}}(\mathrm{X};\mathrm{E})\) , the function \(\mathbf{u} \circ \mathbf{f}\) belongs to \(\widetilde{\mathcal{K}}(\mathrm{X};\mathrm{F})\) and

\[
\int \mathbf {u} (\mathbf {f} (x)) d \mu (x) = ^ {t t} \mathbf {u} \left(\int \mathbf {f} (x) d \mu (x)\right).\tag{2}
\]

For every continuous linear form \(z^{\prime} \in F^{\prime}\) , we have \(z^{\prime} \circ u \circ f = y^{\prime} \circ f\) on setting \(y^{\prime} = z^{\prime} \circ u = {}^{t}u(z^{\prime}) \in E^{\prime}\) , whence the first assertion; moreover, for all \(z^{\prime} \in F^{\prime}\) ,

\[
\begin{array}{r l} \left\langle \int (\mathbf {u} \circ \mathbf {f}) d \mu , \mathbf {z} ^ {\prime} \right\rangle & = \int \langle \mathbf {u} \circ \mathbf {f}, \mathbf {z} ^ {\prime} \rangle d \mu = \int \langle \mathbf {f}, ^ {t} \mathbf {u} (\mathbf {z} ^ {\prime}) \rangle d \mu \\ & = \left\langle \int \mathbf {f} d \mu , ^ {t} \mathbf {u} (\mathbf {z} ^ {\prime}) \right\rangle = \left\langle {} ^ {t t} \mathbf {u} \left(\int \mathbf {f} d \mu\right), \mathbf {z} ^ {\prime} \right\rangle , \end{array}
\]

whence the formula (2).

PROPOSITION 3. --- For every function \(g \in \mathcal{C}(\mathrm{X};\mathbf{C})\) and every function \(\mathbf{f} \in \widetilde{\mathcal{K}}(\mathrm{X};\mathrm{E})\) , the function \(g\mathbf{f}\) belongs to \(\widetilde{\mathcal{K}}(\mathrm{X};\mathrm{E})\) and

\[
\int \mathbf {f} d (g \cdot \mu) = \int \mathbf {f} g d \mu .\tag{3}
\]

Indeed, for every \(z^{\prime} \in E^{\prime}\) , \(\langle gf, z^{\prime} \rangle = g\langle f, z^{\prime} \rangle\) , whence the first assertion; moreover,

\[
\begin{array}{r l} \left\langle \int \mathbf {f} d (g \cdot \mu), \mathbf {z} ^ {\prime} \right\rangle & = \int \langle \mathbf {f}, \mathbf {z} ^ {\prime} \rangle d (g \cdot \mu) = \int g \langle \mathbf {f}, \mathbf {z} ^ {\prime} \rangle d \mu \\ & = \int \langle g \mathbf {f}, \mathbf {z} ^ {\prime} \rangle d \mu = \left\langle \int g \mathbf {f} d \mu , \mathbf {z} ^ {\prime} \right\rangle , \end{array}
\]

whence (3).

PROPOSITION 4. --- Let \(\mu\) be a positive measure on X, S its support, and f a function in \(\widetilde{\mathcal{K}}(\mathrm{X};\mathrm{E})\) . Suppose E is Hausdorff, and equip the space \(\mathrm{E}^{\prime*}\) with the weak topology \(\sigma(\mathrm{E}^{\prime*},\mathrm{E}^{\prime})\) .

\begin{enumerate}
\def\labelenumi{(\roman{enumi})}
\item
  The integral \(\int \mathbf{f}d\mu\) belongs to the closure C in \(\mathrm{E}^{\prime *}\) of the convex cone generated by \(\mathbf{f}(\mathrm{S})\) .
\item
  If \(\mu\) is bounded, then the integral \(\int \mathbf{f} d\mu\) belongs to the set \(\| \mu \| \cdot D\) , where \(D\) is the closed convex envelope of \(\mathbf{f}(S)\) in \(\mathrm{E}'^*\) .
\end{enumerate}

If \(\mathrm{E}\) is complex, we equip \(\mathrm{E}\) with its underlying real vector space structure, which, as we have seen, does not modify the formula (1).

\begin{enumerate}
\def\labelenumi{(\roman{enumi})}
\tightlist
\item
  We know that C is the intersection of the closed half-spaces in \(E^{\prime*}\) that contain \(\mathbf{f}(\mathrm{S})\) and are determined by closed hyperplanes passing through 0; it therefore suffices to prove that, for \(z^{\prime} \in E^{\prime}\) , the relation \(\langle\mathbf{f}(x), \mathbf{z}^{\prime}\rangle \geqslant 0\) for all \(x \in S\) implies
\end{enumerate}

\[
\left\langle \int \mathbf {f} d \mu , \mathbf {z} ^ {\prime} \right\rangle \geqslant 0;
\]

but since

\[
\left\langle \int \mathbf {f} d \mu , \mathbf {z} ^ {\prime} \right\rangle = \int \left\langle \mathbf {f}, \mathbf {z} ^ {\prime} \right\rangle d \mu ,
\]

this follows from §2, No.~3, Cor. 2 of Prop. 8.

\begin{enumerate}
\def\labelenumi{(\roman{enumi})}
\setcounter{enumi}{1}
\tightlist
\item
  We know that D is the intersection of the closed half-spaces in \(E^{\prime*}\) that contain \(\mathbf{f}(S)\) ; it therefore suffices to show that, for \(z^{\prime} \in E^{\prime}\) , the relation \(\langle\mathbf{f}(x), \mathbf{z}^{\prime}\rangle \leqslant a\) for all \(x \in S\) implies
\end{enumerate}

\[
\left\langle \int \mathbf {f} d \mu , \mathbf {z} ^ {\prime} \right\rangle \leqslant a \| \mu \|;
\]

but this follows from §2, No.~3, Cor. 3 of Prop. 8.

COROLLARY. --- Let \(\mu\) be a positive measure on X, f a mapping belonging to \(\mathcal{K}(\mathrm{X};\mathrm{E})\) , and D the closed convex envelope of \(\mathbf{f}(\mathrm{X})\) in \(\mathrm{E}'^*\) . There exists a number \(a > 0\) such that \(\int \mathbf{f} d\mu \in a \cdot \mathrm{D}\) .

If \(\nu\) is any measure on \(X\) , there exist numbers \(a_1, a_2, a_3, a_4 > 0\) such that \(\int \mathbf{f} d\nu \in a_1\mathrm{D} - a_2\mathrm{D} + ia_3\mathrm{D} - ia_4\mathrm{D}\) .

Suppose first that \(\mu\) is positive; by hypothesis, the support K of f is compact; if \(\nu\) is the restriction of \(\mu\) to a relatively compact open neighborhood of K, then \(\nu\) is bounded and \(\int f d\mu = \int f d\nu \in \|\nu\| \cdot D\) by Prop. 4, (ii). The second result follows from this, since any complex measure may be written as \(\mu_{1} - \mu_{2} + i\mu_{3} - i\mu_{4}\) , where the \(\mu_{j}\) are positive.

PROPOSITION 5. --- Suppose that the space X is compact, and let f be a continuous mapping of X into a Hausdorff locally convex space E. The closed convex envelope of \(\mathbf{f}(\mathbf{X})\) in \(E^{\prime*}\) (for \(\sigma(E^{\prime*}, E^{\prime})\) ) is equal to the set of vectors \(\int f d\mu\) for all of the positive measures \(\mu\) on X of total mass 1.

Let C be the closed convex envelope of \(\mathbf{f}(\mathbf{X})\) in \(E^{\prime*}\) ; since \(\mathbf{f}(\mathbf{X})\) is compact and \(E^{\prime*}\) is complete, C is compact. We already know (Prop. 4)

that \(\int \mathbf{f}d\mu \in \mathbb{C}\) for every measure \(\mu\) belonging to the convex set H of positive measures on X of total mass equal to 1. On the other hand, H is convex and compact for the vague topology (§1, No.~9, Cor. 3 of Prop. 15) and is the closure (for this topology) of the convex set \(\mathrm{H}_0\) of positive measures of mass 1 and finite support (§2, No.~4, Cor. 3 of Th. 1). Now, the image of \(\mathrm{H}_0\) under the mapping \(\mu \mapsto \int \mathbf{f}d\mu\) is the convex envelope \(\mathrm{C}_0\) of \(\mathbf{f}(\mathbf{X})\) in \(\mathrm{E}'^*\) . On the other hand, this mapping is continuous for the vague topology on \(\mathcal{M}(\mathrm{X};\mathbf{C})\) and the topology \(\sigma (\mathrm{E}^{\prime *},\mathrm{E}^{\prime})\) on \(\mathrm{E}'^*\) since \(\langle \int \mathbf{f}d\mu ,\mathbf{z}'\rangle = \int \langle \mathbf{f},\mathbf{z}'\rangle d\mu\) by definition; thus the image of \(\mathrm{H} = \overline{\mathrm{H}}_0\) is a compact convex set containing \(\mathrm{C}_0\) and contained in C; since \(\mathrm{C} = \overline{\mathrm{C}}_0\) , this image is equal to C.

PROPOSITION 6. --- For every continuous mapping with compact support f of X into a Hausdorff locally convex space E, every continuous semi-norm q on E and every measure \(\mu\) on X such that \(\int f d\mu \in E\) ,

\[
q \left(\int \mathbf {f} d \mu\right) \leqslant \int (q \circ \mathbf {f}) d | \mu |.\tag{4}
\]

Let D be the set of \(z \in E\) such that \(q(\mathbf{z}) \leqslant 1\) ; D is closed, convex and contains 0, therefore \(D = D^{\circ\circ}\) (TVS, II, §6, No.~3, Cor. 3 of Th. 1). It therefore suffices to prove that for every \(z' \in D^{\circ}\) ,

\[
\left| \left\langle \int \mathbf {f} d \mu , \mathbf {z} ^ {\prime} \right\rangle \right| \leqslant \int (q \circ \mathbf {f}) d | \mu |;
\]

since

\[
\left\langle \int \mathbf {f} d \mu , \mathbf {z} ^ {\prime} \right\rangle = \int \left\langle \mathbf {f}, \mathbf {z} ^ {\prime} \right\rangle d \mu ,
\]

and since, by the definition of \(D^{\circ}\) ,

\[
| \langle \mathbf {f} (x), \mathbf {z} ^ {\prime} \rangle | \leqslant q (\mathbf {f} (x))
\]

for all \(x \in \mathbf{X}\) , the desired inequality follows from the inequality (13) of §1, No.~6.

